\documentclass[11pt]{article}
\usepackage{amsmath,amssymb,amsthm}
\newtheorem{proposition}{Proposition}
\newtheorem{theorem}[proposition]{Theorem}
\newtheorem{lemma}[proposition]{Lemma}
\theoremstyle{remark}
\newtheorem{remark}[proposition]{Remark}
\newtheorem{counterexample}[proposition]{Counterexample}
\title{VAU: Adam-metric correction of the actual AdamW step for LoRA concept erasure,\\ with fresh-batch certified acceptance --- derivation}
\author{XXResearch flow\_method\_design\_v1 (coordinator synthesis; Props.~1 and~6 from Agent C)}
\date{2026-10-19}
\begin{document}
\maketitle

\paragraph{Status.} Every ingredient below is standard (Section~7 lists them). What is offered is (i) a precise
statement of which retention guarantees survive the \emph{actual} AdamW update used by native GEM (LoRA, no retention
term; \texttt{agents/A\_gem\_code\_facts.md} \S5--6), (ii) a strict first-order separation between the Adam-metric correction
and its closest published primitive (Euclidean projection of the Adam update, Adam-NSCL-style; UNVERIFIED source), and (iii) a
finite-step drift bound for one corrected step that makes the sketch, ridge, staleness and curvature errors explicit.
The results are about a linearised retained map plus explicitly bounded remainders; they say nothing about image
semantics, object removal, or FLUX.

\section{Notation and exact objects}
$\theta\in\mathbb{R}^P$: trainable LoRA parameters (all $A,B$ factors flattened), base weights frozen, $\theta_0$ the
initialisation ($B=0$). $L_e(\theta)$: the native GEM erasure loss (detached teachers), $g=\nabla L_e(\theta)$.
Retained map $f_r:\mathbb{R}^P\to\mathbb{R}^{n}$ stacks the student outputs on a retained probe batch $R_1$ ($n=N_r d$);
$J=\partial f_r(\theta)\in\mathbb{R}^{n\times P}$ (never formed densely at scale).
AdamW (torch semantics, $\beta_1,\beta_2,\epsilon,\lambda_{wd}$, step $t$):
$m'=\beta_1 m+(1-\beta_1)g$, $v'=\beta_2 v+(1-\beta_2)g^{\odot2}$, $\hat m=m'/(1-\beta_1^t)$, $\hat v=v'/(1-\beta_2^t)$,
\[ D=\operatorname{diag}(d),\quad d_i=\frac{1}{\sqrt{\hat v_i}+\epsilon}>0,\qquad
\Delta_0=-\eta\,D\hat m-\eta\lambda_{wd}\theta\quad(\text{the actual step}). \]
Sketch: $\Sigma\in\mathbb{R}^{k\times n}$ with i.i.d.\ $\mathcal N(0,1/k)$ entries, $S=\Sigma J\in\mathbb{R}^{k\times P}$
(computed with $k$ VJPs). Metric $H=D^{-1}$, $\|x\|_H^2=x^\top D^{-1}x$.
Candidate correction (exact, with ridge $\mu\ge0$):
\[ \Delta_1=\Delta_0-DS^\top (SDS^\top+\mu I)^{-1}S\Delta_0 . \tag{1}\]
For $\mu=0$ and $S$ of full row rank write $\Pi_D=I-DS^\top(SDS^\top)^{-1}S$; Euclidean analogue
$\Pi=I-S^\top(SS^\top)^{-1}S$. Pre-optimizer projection (GPM/OGD-style) feeds $\Pi g$ into AdamW instead of $g$.
\emph{Idealised} statements set $\mu=0$, $\Sigma=I$ (exact $J$) and fresh $J$; Section~5 removes each idealisation.

\section{Pre-optimizer projection is not retention-neutral under AdamW (Agent C, T1)}
\begin{proposition}\label{p:pre}
Fix $J$ and suppose every gradient fed to AdamW lies in $\mathcal N=\operatorname{null}(J)$ (pre-projection). Then the
momentum $\hat m\in\mathcal N$ and the first-order retained change of the step $\Delta_0$ is
$J\Delta_0=-\eta JD\hat m-\eta\lambda_{wd}J\theta$. (a) $JD\hat m=0$ for every $\hat m\in\mathcal N$ iff
$D\mathcal N\subseteq\mathcal N$. (b) For every $c\in\mathbb{R}$, $\|JD\hat m\|\le\|J\|\,\|D-cI\|\,\|\hat m\|$; with
$c=(d_{\max}+d_{\min})/2$ this gives $\|JD\hat m\|\le\|J\|\frac{d_{\max}-d_{\min}}{2}\|\hat m\|$. (c) The weight-decay
term $\eta\lambda_{wd}J\theta$ is untouched by any gradient projection.
\end{proposition}
\begin{proof}
$\hat m$ is a linear combination of past (projected) gradients, all in the subspace $\mathcal N$, so $\hat m\in\mathcal N$.
(a) $JD\hat m=0\ \forall \hat m\in\mathcal N$ iff $D\mathcal N\subseteq\operatorname{null}(J)=\mathcal N$.
(b) $J\hat m=0$ gives $JD\hat m=J(D-cI)\hat m$; take norms; $\|D-cI\|=\max_i|d_i-c|$, minimised at the midrange.
(c) The term does not depend on $g$.
\end{proof}
\begin{remark} Since $D$ is diagonal with (generically) distinct entries, $D$-invariant subspaces are sums of coordinate
subspaces; a null space that is \emph{rotated} relative to the coordinate axes is generically not invariant. At AdamW step
1, $\hat m/\sqrt{\hat v}=\operatorname{sign}(g)$ entrywise ($|g_i|\gg\epsilon$), so the retained change is
$\eta\|J\operatorname{sign}(\Pi g)\|$, which does not shrink as the projection becomes more accurate. \end{remark}

\section{The Adam-metric correction: characterisation and a strict separation}
\begin{proposition}[trust-region characterisation; standard]\label{p:tr}
Let $D\succ0$, $S$ full row rank, $r>0$, and $u=D\hat m$ with $\hat m=g$. The problem
\[ \max_x\ -g^\top x\quad\text{s.t.}\quad Sx=0,\ \|x\|_H\le r \]
has the unique solution $x^\star=-r\,\Pi_D Dg/\|\Pi_D Dg\|_H$ whenever $\Pi_D Dg\ne0$, and the optimal value is
$r\|\Pi_D Dg\|_H$. Moreover
\[ -g^\top(-\Pi_D Dg)=\|(I-Q)D^{1/2}g\|^2\ \ge 0,\qquad Q=D^{1/2}S^\top(SDS^\top)^{-1}SD^{1/2}, \tag{2}\]
with equality iff $g\in\operatorname{range}(S^\top)$. Also $\|\Pi_D w\|_H\le\|w\|_H$ for all $w$, and $S\Pi_D=0$.
\end{proposition}
\begin{proof}
Substitute $x=D^{1/2}y$: maximise $-(D^{1/2}g)^\top y$ s.t.\ $(SD^{1/2})y=0$, $\|y\|\le r$. The feasible set is the ball
intersected with $\operatorname{null}(SD^{1/2})$, whose orthogonal projector is $I-Q$ ($Q$ is the orthogonal projector onto
$\operatorname{range}(D^{1/2}S^\top)$). By Cauchy--Schwarz on that subspace the unique maximiser is
$y^\star=-r(I-Q)D^{1/2}g/\|(I-Q)D^{1/2}g\|$, and $D^{1/2}(I-Q)D^{1/2}g=\Pi_D Dg$, giving $x^\star$. Then
$g^\top\Pi_D Dg=(D^{1/2}g)^\top(I-Q)(D^{1/2}g)=\|(I-Q)D^{1/2}g\|^2$ since $I-Q$ is an orthogonal projector. Equality
iff $D^{1/2}g\in\operatorname{range}(D^{1/2}S^\top)$ iff $g\in\operatorname{range}(S^\top)$. In $y$-coordinates
$\Pi_D$ acts as the orthogonal projector $I-Q$, hence non-expansive in $\|\cdot\|_H$; $S\Pi_D=S-SDS^\top(SDS^\top)^{-1}S=0$.
\end{proof}

\begin{proposition}[strict separation from Euclidean post-projection]\label{p:sep}
For the Euclidean correction $\Delta_E=\Pi\Delta_0$ of the same AdamW step ($\hat m=g$, $\lambda_{wd}=0$), the first-order
erasure progress $-g^\top\Delta_E=\eta\,g^\top\Pi Dg$ \emph{can be negative}; the Adam-metric progress
$-g^\top\Pi_D\Delta_0=\eta\,g^\top\Pi_DDg$ never is (Prop.~\ref{p:tr}). The two coincide when $D=cI$.
Explicit instance: $P=2$, $S=[1,1]$, $D=\operatorname{diag}(0.01,1)$, $g=(1,0.5)$:
$g^\top\Pi Dg=\tfrac12(d_1g_1-d_2g_2)(g_1-g_2)=\tfrac12(-0.49)(0.5)=-0.1225<0$, while
$g^\top\Pi_DDg=\|(I-Q)D^{1/2}g\|^2=\frac{d_1d_2(g_1-g_2)^2}{d_1+d_2}=\frac{0.01\cdot0.25}{1.01}\approx0.002475>0$.
\end{proposition}
\begin{proof}
For $S=[1,1]$, $\Pi=\frac12\begin{psmallmatrix}1&-1\\-1&1\end{psmallmatrix}$, so $\Pi Dg=\frac12(d_1g_1-d_2g_2)(1,-1)^\top$ and
$g^\top\Pi Dg=\frac12(d_1g_1-d_2g_2)(g_1-g_2)$. For the metric version, $SDS^\top=d_1+d_2$ and
$\Pi_DDg=Dg-\frac{d_1g_1+d_2g_2}{d_1+d_2}(d_1,d_2)^\top=\frac{d_1d_2(g_1-g_2)}{d_1+d_2}(1,-1)^\top$, so
$g^\top\Pi_DDg=\frac{d_1d_2(g_1-g_2)^2}{d_1+d_2}$. Substitute. If $D=cI$ then $\Pi_D=\Pi$.
\end{proof}
\begin{remark}[what Prop.~\ref{p:sep} does \emph{not} say] It does not say the metric correction always yields more
progress: with $g=(1,-0.5)$ in the same instance the Euclidean progress is larger (0.3825 vs 0.0223 in units of $\eta$;
\texttt{code/scratch\_checks/sign\_check.py}). Prop.~\ref{p:tr} gives optimality only at equal $H$-norm (Adam's own trust
region). Comparisons must therefore be made at tuned or matched step size --- the protocol does so.
Under random draws (development only) the Euclidean sign failure occurred in 3.5\,\% of cases (same script).\end{remark}

\begin{counterexample}[momentum]\label{c:mom}
With $\hat m\ne g$, $-g^\top\Pi_D D\hat m$ can be negative (take $\hat m$ dominated by an old gradient opposite to $g$
outside $\operatorname{range}(S^\top)$). Hence the rule ``if $g^\top\Delta_1\ge0$ set $\Delta_1=0$''; Prop.~\ref{p:tr}'s
sign guarantee covers only $\beta_1=0$ or $\hat m\parallel g$ on $\operatorname{null}(S)$-relevant directions.
\end{counterexample}

\begin{remark}[weight decay] Because (1) acts on the whole $\Delta_0$, the first-order retained change of the decoupled
weight-decay term is also removed ($S\Delta_1=0$ exactly for $\mu=0$), unlike Prop.~\ref{p:pre}(c). For LoRA the
weight-decay map change is exactly $((1-\eta\lambda)^2-1)BA$ (both factors scale), whose second-order part is not removed.
\end{remark}

\section{Exact finite-step behaviour of one corrected step}
Let $\theta^+=\theta+s\Delta_1$. The exact retained change is $f_r(\theta^+)-f_r(\theta)$; write the full (unsketched)
Jacobian at the current point $J_\theta$ and at the probe refresh point $J_{\theta_s}$ (staleness).
Assumptions: (A1) $\|J_\theta-J_{\theta_s}\|\le L_J\|\theta-\theta_s\|$; (A2) along the segment $[\theta,\theta^+]$ the
second derivative of $f_r$ is bounded: $\|f_r(\theta+h)-f_r(\theta)-J_\theta h\|\le\frac{\beta}{2}\|h\|^2$.

\begin{lemma}[ridge residual]\label{l:ridge}
With $S$ any $k\times P$ matrix, (1) gives $S\Delta_1=\mu(SDS^\top+\mu I)^{-1}S\Delta_0$, hence
$\|S\Delta_1\|\le\frac{\mu}{\lambda_{\min}(SDS^\top)+\mu}\|S\Delta_0\|$.
\end{lemma}
\begin{proof} $S\Delta_1=S\Delta_0-SDS^\top(G+\mu I)^{-1}S\Delta_0$ with $G=SDS^\top$; $I-G(G+\mu I)^{-1}=\mu(G+\mu I)^{-1}$;
the eigenvalues of $\mu(G+\mu I)^{-1}$ are $\mu/(\lambda_i+\mu)$. \end{proof}

\begin{lemma}[sketch residual]\label{l:sketch}
Let $J=J_{\theta_s}$ have singular values $\sigma_1\ge\sigma_2\ge\dots$, $\Sigma$ Gaussian $k\times n$, and $Y=\operatorname{range}(J^\top\Sigma^\top)\subseteq\mathbb{R}^P$. Every $x$ with $\Sigma Jx=0$ satisfies $x\perp Y$ and
$\|Jx\|\le E_\Sigma\|x\|$ with $E_\Sigma=\|J(I-P_Y)\|$. (a) If $k\ge\operatorname{rank}J$ then $E_\Sigma=0$ almost surely.
(b) For $k=k_0+p$, $p\ge2$, the randomized range-finder bound (Halko--Martinsson--Tropp 2011, \S10, average spectral error
theorem; applied to $A=J^\top$, $\Omega=\Sigma^\top$) gives
$\mathbb{E}E_\Sigma\le\bigl(1+\sqrt{k_0/(p-1)}\bigr)\sigma_{k_0+1}+\frac{e\sqrt{k}}{p}\bigl(\sum_{j>k_0}\sigma_j^2\bigr)^{1/2}$,
so by Markov $\Pr(E_\Sigma>t\,\mathbb{E}E_\Sigma)\le1/t$.
\end{lemma}
\begin{proof}
$\Sigma Jx=0\iff x\perp\operatorname{range}(J^\top\Sigma^\top)=Y$, so $x=(I-P_Y)x$ and $\|Jx\|=\|J(I-P_Y)x\|\le E_\Sigma\|x\|$.
(a) A Gaussian $\Sigma^\top$ restricted to $\operatorname{range}(J)$ (dimension $\operatorname{rank}J\le k$) is a.s.\ of full
rank $\operatorname{rank}J$, so $Y=\operatorname{range}(J^\top)$ and $J(I-P_Y)=0$. (b) $\|J(I-P_Y)\|=\|(I-P_Y)J^\top\|$, which is
exactly the quantity bounded in HMT \S10 for $A=J^\top$ with Gaussian test matrix $\Omega=\Sigma^\top$ (scaling of $\Omega$ does
not change $Y$). The constant is quoted from that theorem from memory (UNVERIFIED theorem number in this sprint; arXiv
unreachable); Markov's inequality gives the tail. Note: the bound requires $\Sigma$ independent of $J$, which holds
because $\Sigma$ is drawn fresh at each refresh before $J$-products are taken; it does \emph{not} require independence from
$x$, because the bound is uniform over $x\perp Y$.
\end{proof}

\begin{theorem}[finite-step retained drift of one VAU step]\label{t:step}
Under (A1)--(A2), with $S=\Sigma J_{\theta_s}$ and $\Delta_1$ from (1),
\[ \|f_r(\theta+s\Delta_1)-f_r(\theta)\|\ \le\ s\Bigl(E_\Sigma\|\Delta_1\|+\kappa_\Sigma\|S\Delta_1\|\Bigr)
 + s\,L_J\|\theta-\theta_s\|\,\|\Delta_1\| + \tfrac{\beta}{2}s^2\|\Delta_1\|^2 , \]
where $\|S\Delta_1\|$ is bounded by Lemma~\ref{l:ridge} (zero if $\mu=0$), $E_\Sigma$ by Lemma~\ref{l:sketch} (zero a.s.\ if
$k\ge\operatorname{rank}J_{\theta_s}$), and $\kappa_\Sigma=1/\sigma_{\min}(\Sigma|_V)$ with $V=J_{\theta_s}Y$, finite (shown in the proof).
\end{theorem}
\begin{proof}
Write $f_r(\theta+h)-f_r(\theta)=J_\theta h+\rho$ with $\|\rho\|\le\frac\beta2\|h\|^2$ (A2), $h=s\Delta_1$. Split
$J_\theta h=J_{\theta_s}h+(J_\theta-J_{\theta_s})h$; the second term is bounded by (A1). Decompose $h=h_\perp+h_Y$ with $h_\perp\perp Y$, $h_Y\in Y$ ($\|h_\perp\|,\|h_Y\|\le\|h\|$). Then
$\|J_{\theta_s}h_\perp\|\le E_\Sigma\|h\|$ (Lemma~\ref{l:sketch}) and $Sh=\Sigma J_{\theta_s}h_Y$ because $\Sigma J_{\theta_s}h_\perp=0$.
$J_{\theta_s}h_Y\in V=J_{\theta_s}Y=\operatorname{range}(J_{\theta_s}J_{\theta_s}^\top\Sigma^\top)$. $\Sigma$ is injective on $V$:
if $z=J_{\theta_s}J_{\theta_s}^\top\Sigma^\top w$ and $\Sigma z=0$ then $w^\top\Sigma z=\|J_{\theta_s}^\top\Sigma^\top w\|^2=0$, so $z=0$.
Hence $\sigma_{\min}(\Sigma|_V)>0$ and $\|J_{\theta_s}h_Y\|\le\kappa_\Sigma\|Sh\|$. If $\mu=0$ then $Sh=0$. Combine with the
triangle inequality. \end{proof}
\begin{remark} Theorem~\ref{t:step} is a \emph{prediction-side} bound: every constant ($E_\Sigma$, $L_J$, $\beta$) is
unknown for a real network. That is why the candidate does not rely on it for acceptance; it relies on Theorem~\ref{t:gate}.\end{remark}

\begin{counterexample}[LoRA gauge; outside the first-order guarantee]\label{c:gauge}
For a LoRA layer $W+BA$ and any $K$, the step $dB=BK$, $dA=-KA$ has first-order map change $dB\,A+B\,dA=0$ for every input,
so it passes any exact first-order constraint, yet the exact change is $dB\,dA=-BK^2A\ne0$ (Agent C, C-X2; computed
2.3\,\% of $\|BAx\|$ at $\|K\|=0.3$). Only an evaluation of the actual finite step detects it.
\end{counterexample}

\section{Certified acceptance under adaptivity, with optimizer-state rollback}
Clipped per-sample drift $r_x(\theta)=\min(\|f_x(\theta)-f_x(\theta_0)\|^2,M)\in[0,M]$; population
$\rho_M(\theta)=\mathbb{E}_{x\sim\mathcal R}r_x(\theta)$ over the retained distribution $\mathcal R$.
\begin{theorem}[gate]\label{t:gate}
Run any update rule for $T$ steps with at most $R$ trials per step. At each trial $j$ draw a fresh batch $R_2^{(j)}$ of
$n_2$ i.i.d.\ samples from $\mathcal R$, independent of everything used to construct the trial point $\vartheta_j$
(conditionally on the history), and accept iff $\hat\rho_j+c\le\tau^2$, $\hat\rho_j=\frac1{n_2}\sum_{x\in R_2^{(j)}}r_x(\vartheta_j)$,
$c=M\sqrt{\ln(TR/\delta)/(2n_2)}$. Then with probability at least $1-\delta$ every accepted iterate satisfies
$\rho_M\le\tau^2$. If additionally $\|f_x(\theta)-f_x(\theta_0)\|^2\le M$ a.s.\ along the run, $\rho_M$ is the true mean drift.
If on rejection $(\theta,m,v,t)$ are restored, the accepted sequence is exactly an AdamW trajectory in which rejected
steps are skipped (no optimizer-state contamination by rejected trials).
\end{theorem}
\begin{proof}
Let $\mathcal F_{j}$ be the $\sigma$-field of everything before batch $j$ is drawn; $\vartheta_j$ is $\mathcal F_j$-measurable
and $R_2^{(j)}$ is independent of $\mathcal F_j$. Conditionally on $\mathcal F_j$, the $r_x(\vartheta_j)$ are i.i.d.\ in $[0,M]$ with
mean $\rho_M(\vartheta_j)$, so Hoeffding gives $\Pr(\rho_M(\vartheta_j)>\hat\rho_j+c\mid\mathcal F_j)\le\delta/(TR)$. The bad
event ``some accepted $\vartheta_j$ has $\rho_M>\tau^2$'' is contained in $\bigcup_j\{\rho_M(\vartheta_j)>\hat\rho_j+c\}$ (acceptance
means $\hat\rho_j+c\le\tau^2$); at most $TR$ trials, so the union bound gives $\le\delta$. The clipping statement is immediate.
Rollback: the accepted update applies the tentative state $(m',v',t+1)$ computed from the same $g$ that produced
$\Delta_0$; a rejected trial leaves $(\theta,m,v,t)$ bitwise unchanged, so the state after each step equals that of AdamW
run on the subsequence of accepted steps (with the projected/scaled displacement in place of the raw step).
\end{proof}
\begin{counterexample}[adaptivity; why $R_1\ne R_2$] If the trial point is constructed in the null space of the same
batch used to verify it, the empirical drift can be zero while population drift is not (Agent C computed $3\cdot10^{-30}$
vs $1.1\cdot10^{-3}$). Theorem~\ref{t:gate} requires independence.\end{counterexample}
\begin{counterexample}[rare subgroup] The gate certifies a mean. A retained subgroup of mass $q$ is absent from
$n_2$ samples with probability $(1-q)^{n_2}$ ($0.37$ for $q=0.005$, $n_2=200$); its drift is uncontrolled.\end{counterexample}

\section{Limiting cases}
(i) $D=cI$ (e.g.\ SGD, or AdamW when all $|g_i|$ equal): metric, Euclidean post- and (for $\beta_1=0$, $\lambda_{wd}=0$)
pre-projection coincide. (ii) $S=0$ (no probes): $\Delta_1=\Delta_0$, VAU-noGate is AdamW. (iii) $g\in\operatorname{range}(S^\top)$:
$\Delta_1$ has zero first-order progress (Prop.~\ref{p:tr}); rejected. (iv) Rich retained set, $\operatorname{rank}J=P$: the
exact null space is $\{0\}$, $\Delta_1\approx0$ (ridge only) --- VAU erases nothing; a soft penalty or the gate alone is
then the only meaningful arm. (v) $\tau\to\infty$: the gate always accepts. (vi) $\tau\to0$: only steps with zero
clipped empirical drift and $c\le\tau^2$ pass; for finite $n_2$, $c>0$, so nothing passes.

\section{Standard ingredients versus what is specific here}
Standard: scaled/metric gradient projection and its trust-region form (Prop.~\ref{p:tr}; e.g.\ Bertsekas' scaled gradient
projection); null-space projection of an optimizer's update for continual learning (Adam-NSCL; UNVERIFIED this sprint);
GPM/OGD gradient projection; HMT randomized range finder (Lemma~\ref{l:sketch}); Taylor/Lipschitz remainders
(Theorem~\ref{t:step}); Hoeffding with union bound and backtracking (Theorem~\ref{t:gate}; TRPO/CPO-style line search).
Specific to this sprint (not claimed as a general contribution): the identification that native GEM's verified training
configuration (AdamW, zero-init LoRA, no retention) sits in the regime of Prop.~\ref{p:pre}; the explicit sign
separation of Prop.~\ref{p:sep} between the metric and Euclidean corrections of the \emph{same} Adam step; the
combination of those with Theorem~\ref{t:gate}'s rollback semantics; and the prospective, falsifiable predictions tested
in \texttt{results/}. Whether any of this improves real erasure/preservation is unknown.
\end{document}
